\chapter{Electric Shock Sensation (Formication)}

\section*{Definition}
A sudden, brief sensation resembling an electric shock, tingling, or crawling feeling on or under the skin, typically lasting seconds, and often occurring without an external stimulus.

\section*{What Happens in Your Body}
Electric shock sensations, medically termed paresthesias, arise from ectopic impulse generation in peripheral sensory nerves or central sensory pathways. Oestrogen supports myelin integrity and nerve conduction; declining oestrogen during perimenopause may increase nerve irritability. Reduced progesterone (which has a membrane-stabilising effect) further lowers the nerve excitation threshold. Hormonal fluctuations alter ion channel function (sodium, potassium, calcium) in sensory neurons, leading to spontaneous depolarisation.

\section*{Causes}
\begin{itemize}
  \item Perimenopausal and menopausal hormonal fluctuations
  \item Cervical or lumbar nerve root compression (radiculopathy)
  \item Peripheral neuropathy (diabetic, alcoholic, nutritional)
  \item Multiple sclerosis (central demyelination)
  \item Migraine aura (sensory variant)
  \item Anxiety and hyperventilation syndrome
  \item Medication side effects (antidepressants, anticonvulsants)
  \item Vitamin B12 deficiency
  \item Hypocalcaemia or hypomagnesaemia
  \item Transient ischaemic attack (sensory cortex involvement)
\end{itemize}

\section*{When to See a Doctor}
See your doctor if:
\begin{itemize}
  \item Symptoms are severe or getting worse over time
  \item Frequent or persistent electric shock sensations with weakness, numbness, vision changes, or other neurological symptoms
  \item You have other concerning symptoms such as fever, weight loss, or bleeding
\end{itemize}

\section*{Related Symptoms}
\begin{itemize}
  \item Tingling numbness
  \item Nerve pain
  \item Dizziness
  \item Anxiety
  \item Headache
\end{itemize}

\textbf{Important:} If this symptom persists, your doctor will check for serious underlying causes.

\section*{Self-Care / Home Management}
\begin{itemize}
  \item Rest and avoid activities that make it worse.
  \item Maintain adequate hydration and a balanced diet.
  \item Over-the-counter remedies may provide symptomatic relief (consult a pharmacist or doctor).
  \item Monitor symptoms and seek professional advice if they persist or worsen.
  \item Avoid self-medicating with prescription medications without medical guidance.
\end{itemize}

\section*{How Your Doctor Will Evaluate You}
\begin{itemize}
  \item A thorough history and physical examination are the first steps in evaluation.
  \item Laboratory tests (blood work, urinalysis) may be ordered based on clinical suspicion.
  \item Imaging studies (X-ray, ultrasound, CT, MRI) may be indicated when structural pathology is suspected.
  \item Specialist referral is arranged when the diagnosis remains uncertain or requires specific expertise.
\end{itemize}

\section*{Treatment Options}
\begin{itemize}
  \item Treatment targets the underlying cause identified through diagnostic evaluation.
  \item Supportive care and symptom management are provided while awaiting definitive therapy.
  \item Medications, lifestyle modifications, and physical therapies may be prescribed as appropriate.
  \item Follow-up ensures treatment effectiveness and allows timely adjustments to the care plan.
\end{itemize}

\clearpage